\section{A dependent pair of components}\label{sec:dependent-pair}

The first step is to find two independent constant linear combinations of the components that are algebraically dependent. The following theorem, proved in the remainder of this section, establishes this reduction.

\begin{theorem}[Dependent pair]\label{thm:dependent-pair}
Let $H\in K[x,y,z]^3$ satisfy $H(0)=0$. Suppose that $JH$ is nilpotent and that the components of $H$ are linearly independent over $K$. Then there exist linearly independent forms $\ell_1,\ell_2\in(K^3)^*$ such that $\ell_1(H)$ and $\ell_2(H)$ are algebraically dependent over $K$.
\end{theorem}

We first work over an algebraically closed field $k$. Since $\det JH=0$, the Jacobian rank is at most two. The dependence result recalled in Section~\ref{sec:preliminaries} excludes rank at most one. Hence $\rank JH=2$. After a simultaneous permutation of source and target coordinates, write $H=(u,v,w)$ with $u,v$ algebraically independent. Set
\[
 \mathcal K=k(x,y,z),\qquad E=k(u,v,w),
\]
and let $L$ be the relative algebraic closure of $E$ in $\mathcal K$. Then $L/k(u,v)$ is finite separable. The derivations $\partial_u,\partial_v$ extend uniquely to $L$ and commute. Put
\begin{equation}\label{eq:surface-derivatives}
 p=w_u,\quad q=w_v,\qquad a=w_{uu},\quad b=w_{uv},\quad c=w_{vv}.
\end{equation}

\subsection{The embedded generic curve}
Consider the Jacobian derivation
\[
 D(f)=\det J(u,v,f),\qquad (C_1,C_2,C_3)=\nabla u\times\nabla v.
\]
By the differential criterion for algebraic dependence,
\begin{equation}\label{eq:kernel-D}
 \ker_{\mathcal K}D=L.
\end{equation}
Indeed, $D(f)=0$ is equivalent to $f$ being algebraic over $k(u,v)$, and $E/k(u,v)$ is algebraic.

Since $\dd w=p\,\dd u+q\,\dd v$, the principal-minor identity is
\begin{equation}\label{eq:plane-sigma}
 \sigma_2(JH)=C_3-pC_1-qC_2.
\end{equation}
The right-hand side is $D(z-px-qy)$, because $p,q\in L$. Thus
\begin{equation}\label{eq:plane-s}
 s=z-px-qy\in L.
\end{equation}
Differentiating this identity, and setting
\[
 A=ax+by+s_u,\qquad B=bx+cy+s_v,
\]
we obtain
\begin{equation}\label{eq:plane-differentials}
 Au_x+Bv_x=-p,\qquad Au_y+Bv_y=-q,\qquad Au_z+Bv_z=1.
\end{equation}
These equations imply
\[
 \tr JH=u_x+v_y+pu_z+qv_z=-AC_1-BC_2.
\]
The coefficients of $A,B$, viewed as affine expressions in $x,y$, lie in $L$. It follows that
\begin{equation}\label{eq:quadratic-tau}
 \tau=\tfrac12ax^2+bxy+\tfrac12cy^2+s_ux+s_vy\in L.
\end{equation}

\begin{proposition}\label{prop:line-conic}
Let $\mathcal C=\Spec L[x,y,z]$, where $L[x,y,z]$ denotes the subalgebra of $\mathcal K$ generated by the source coordinates. Then $\mathcal C$ is a geometrically integral embedded affine line or plane conic. In the conic case its points at infinity have directions
\begin{equation}\label{eq:infinity-directions}
 e=(\xi,\eta,p\xi+q\eta),\qquad a\xi^2+2b\xi\eta+c\eta^2=0,
\end{equation}
over $\overline L$.
\end{proposition}

\begin{proof}
This algebra has fraction field $\mathcal K$ and dimension one over $L$. It is geometrically integral by Lemma~\ref{lem:field-facts}. Equations~\eqref{eq:plane-s} and~\eqref{eq:quadratic-tau} place the curve in the plane $Z=pX+qY+s$ and on the quadratic
\[
 \Phi(X,Y)=\tfrac12aX^2+bXY+\tfrac12cY^2+s_uX+s_vY-\tau.
\]
The coefficients $a,b,c$ are not all zero. Otherwise, both coordinate derivations on $L$ annihilate $p,q$, so they belong to $k$. The same argument gives $w-pu-qv\in k$. Evaluating at the origin then contradicts linear independence.

The kernel of $L[X,Y]\to \mathcal K$, $X\mapsto x$, $Y\mapsto y$, is a height-one prime containing the nonzero polynomial $\Phi$. It is generated by an irreducible factor of $\Phi$, of degree one or two. If the degree is two, that factor is a scalar multiple of $\Phi$. Homogenizing the plane and quadratic equations gives~\eqref{eq:infinity-directions}.
\end{proof}

The direction space of the containing affine plane is
\begin{equation}\label{eq:tangent-plane}
 \mathcal T=\{(\xi,\eta,p\xi+q\eta):\xi,\eta\in L\}.
\end{equation}
After extending scalars to $\mathcal K$, this is the column space of $JH$. No nonzero constant linear form vanishes on $\mathcal T$, since such a form would satisfy $\ell JH=0$, and hence $\ell(H)=0$ by normalization.

The construction is compatible with constant conjugation. For $S\in\GL_3(k)$, regard $x'=Sx$ and $H'=SH$ as elements of the same field $\mathcal K$. Then $k(H')=E$, $L[x']=L[x]$, and the Jacobian in the new coordinates is $S(JH)S^{-1}$. The transformation $S$ acts on the embedded curve, its directions, and $\mathcal T$. We may therefore repeat the construction with any algebraically independent pair of the new components.

\subsection{Constant planes at infinity}
Let $H_d\neq0$ be the highest homogeneous part of $H$. Taking the highest homogeneous part of the identity $(JH)^3=0$ gives $(JH_d)^3=0$. By~\eqref{eq:homogeneous-form}, a constant conjugation puts $H_d$ in the form $(0,\gamma x^d,B(x,y))$. If $\gamma\neq0$, its common zero set is contained in $x=0$. If $\gamma=0$, the nonzero binary form $B$ factors into linear forms over $k$. In either case there are finitely many nonzero constant linear forms $m_1,\ldots,m_N$ such that
\begin{equation}\label{eq:finite-planes}
 H_d(e)=0\quad\Longrightarrow\quad\prod_{j=1}^N m_j(e)=0.
\end{equation}
This implication holds over every extension of $k$.

The polynomial $H_i(X,Y,Z)-H_i(x,y,z)$, with coefficients in $L$, vanishes on $\mathcal C$. Homogenizing to the common degree $d$ shows that every point at infinity of its projective closure satisfies $H_d(e)=0$. Therefore every direction at infinity lies in one of the constant planes in~\eqref{eq:finite-planes}.

\subsection{Excluding a line}
Suppose that $\mathcal C$ is a line. Choose a constant form $m$ that vanishes on its direction and use $m$ as the first row of a simultaneous change of coordinates. The new coordinate $x$ is constant on the geometric line, so it is algebraic over $L$. Relative algebraic closedness gives $x\in L$.

Put $k_0=k(x)$. As elements of $k_0[y,z]$, the components of $H$ generate a field of transcendence degree one over $k_0$. Choose a closed generator $r\in k_0[y,z]$. Then
\begin{equation}\label{eq:line-field}
 L=k_0(r).
\end{equation}
Indeed, $r$ is algebraic over $E(x)$ and thus belongs to $L$. Conversely, every element of $L$ is algebraic over $k_0(r)$, which is relatively algebraically closed in $\mathcal K$.

The kernel defining the curve in the $Y,Z$ coordinates is generated by $r(Y,Z)-r$. To see this, first use the prime ideal $(r(Y,Z)-R)$ in $k_0[R,Y,Z]$ and then localize at $k_0[R]\setminus\{0\}$. Because the curve is a line, $r$ has degree one in $Y,Z$. We can write
\[
 r=a(x)y+b(x)z+c(x),\qquad a(x),b(x),c(x)\in k(x).
\]
If $a(x)=0$, all components of $H$ are independent of $y$. Otherwise, replace the generator by $s=y-\rho(x)z$ and write
\[
 H=(f(x,s),g(x,s),h(x,s)),\qquad f,g,h\in k(x)[s].
\]
The trace equation, in the independent variables $x,s,z$, becomes
\[
 f_x+g_s-\rho h_s-\rho' zf_s=0.
\]
Consequently $\rho'f_s=0$. If $\rho'\neq0$, then the first component in the source coordinates belongs to $k(x)\cap k[x,y,z]=k[x]$. Its Jacobian row is $(f'(x),0,0)$, and nilpotence forces $f'=0$. Normalization gives $f=0$, contrary to independence. Thus $\rho\in k$.

In either case, $H$ is invariant under translation in a fixed nonzero source direction. A constant conjugation sends that direction to the $z$-axis, giving a map $(A(x,y),B(x,y),C(x,y))$. The planar Jacobian of $(A,B)$ is nilpotent. The planar dependence result recalled in Section~\ref{sec:preliminaries} gives a constant relation between $A,B$, again contradicting independence. Hence $\mathcal C$ is a conic.

\subsection{Excluding two points at infinity}
Suppose that the conic has two distinct directions $e_1,e_2$ over $\overline L$. Choose constant forms $m_1,m_2$ with $m_i(e_i)=0$. Their restrictions to $\mathcal T$ are nonzero. On $\mathcal T\otimes_L\overline L$, their kernels are the distinct lines spanned by $e_1,e_2$, so the restrictions are linearly independent over $k$. Hence $u=m_1(H)$ and $v=m_2(H)$ are algebraically independent. Use their restrictions to $\mathcal T$ as the first two rows of a simultaneous coordinate change and complete the third row.

The projected directions in the new $x,y$ plane are $[0:1]$ and $[1:0]$. Recomputing~\eqref{eq:infinity-directions} gives $a=c=0$ and $b\neq0$. Commutation of derivatives gives $b_u=b_v=0$, and therefore $b=\gamma\in k^*$. Integration twice in the function field, (3.1) gives
\begin{equation}\label{eq:bilinear-graph}
 w=\gamma uv+\alpha u+\beta v+\delta,\qquad \alpha,\beta,\delta\in k.
\end{equation}
The following lemma excludes this possibility.

\begin{lemma}\label{lem:bilinear-exclusion}
A polynomial map $(u,v,\gamma uv+\alpha u+\beta v+\delta)$ with $\gamma\neq0$ cannot have both a nilpotent Jacobian and Jacobian rank two.
\end{lemma}

\begin{proof}
Suppose first that $m=\deg_z u>0$ and $n=\deg_zv>0$. The trace equation is
\begin{equation}\label{eq:bilinear-trace}
 u_x+v_y+\gamma(uv)_z+\alpha u_z+\beta v_z=0.
\end{equation}
The degree of $(uv)_z$ forces $m+n-1\leq\max(m,n)$, so $\min(m,n)=1$. If $m>n=1$, let $A,B\in k[x,y]\setminus\{0\}$ be the leading coefficients in $z$. The coefficient of $z^m$ is
\[
 A_x+(m+1)\gamma AB=0.
\]
Comparing degrees in $x$ over $k(y)$ gives a contradiction. The case $n>m=1$ is symmetric. If $m=n=1$, we obtain $A_x+B_y+2\gamma AB=0$. The total degree of $AB$ exceeds that of either nonzero derivative, again yielding a contradiction.

After interchanging the first two coordinates if necessary, assume $u_z=0$. Using the trace equation to eliminate $v_y$ from the principal-minor identity gives
\begin{equation}\label{eq:bilinear-minors}
 u_x^2+u_y\bigl(v_x+(\gamma v+\alpha)v_z\bigr)=0.
\end{equation}
If $u_y=0$, then $u_x=0$, so $u$ is constant and the Jacobian rank is at most one. If $v_z=0$, the third column of the Jacobian is zero and its first two rows are dependent by $\sigma_2=0$; again the rank is at most one. Thus $u_yv_z\neq0$.

If $\deg_zv\geq2$, the term $vv_z$ in~\eqref{eq:bilinear-minors} has strictly larger $z$-degree than the other terms. Therefore $v=Cz+D$ with $C\neq0$. The coefficient of $z$ gives $C_x+\gamma C^2=0$, which is impossible by comparing degrees in $x$.
\end{proof}

\subsection{The remaining direction}
The conic therefore has a unique point at infinity, counted twice. Let $e$ be its direction and choose a nonzero constant form $m$ with $m(e)=0$. Use $m$ as the first row of a simultaneous coordinate change. The first component $u=m(H)$ is nonconstant. Choose a constant combination $v$ algebraically independent of $u$, complete the coordinate change, and repeat the construction of the curve.

Now $e_x=0$, and~\eqref{eq:tangent-plane} implies $e_y\neq0$. The double root of the binary quadratic in~\eqref{eq:infinity-directions} is $[0:1]$. Hence $a\neq0$ and $b=c=0$. Both coordinate derivatives of $w_v$ vanish, so $w_v=\mu\in k$. It follows that
\[
 \dd(w-\mu v)=w_u\,\dd u.
\]
We give the field-theoretic justification for the resulting algebraic dependence. The extension $L/k(u,v)$ is finite separable, and the extended coordinate derivation $\partial_v$ satisfies
\[
 \partial_v(w-\mu v)=w_v-\mu=0.
\]
By Lemma~\ref{lem:field-facts}(iii), $w-\mu v$ is therefore algebraic over $k(u)$. More explicitly, let $P(T)\in k(u,v)[T]$ be its monic minimal polynomial. Differentiating $P(w-\mu v)=0$ with respect to $v$ gives $(\partial_vP)(w-\mu v)=0$, where $\partial_vP$ denotes coefficientwise differentiation. Since $P$ is monic, $\partial_vP$ is either zero or has degree smaller than $P$. Minimality forces $\partial_vP=0$. The coefficients of $P$ thus belong to the constant field of $\partial_v$ on $k(u,v)$, which is $k(u)$. Clearing their denominators yields a nonzero polynomial relation over $k$ between $u$ and $w-\mu v$. Thus $u$ and $w-\mu v$ are algebraically dependent over $k$. The corresponding constant forms are independent. This proves Theorem~\ref{thm:dependent-pair} over an algebraically closed field.

\subsection{Descent and completion of the proof}
We use Proposition~\ref{prop:pair-normal}, which is proved in the next section over an arbitrary characteristic-zero field, assuming that a dependent pair is already defined over that field. Its proof does not use Theorem~\ref{thm:dependent-pair}.

Extend scalars from $K$ to $\Kbar$. Linear independence of the components is preserved, since it is a rank condition on their coefficient matrix. The preceding argument and Proposition~\ref{prop:pair-normal} give the canonical form over $\Kbar$. Consider the finite-dimensional relation space
\begin{equation}\label{eq:relation-space}
 V_K=\{R\in K[U,V,W]_{\leq2}:R(H_1,H_2,H_3)=0\}.
\end{equation}
It is the kernel of a linear map of finite-dimensional coefficient spaces over $K$, and therefore
\[
 V_K\otimes_K\Kbar=V_{\Kbar}.
\]
For $\Ncal_P$, the first two components are algebraically independent and the third is the square of the first. Its full ideal of relations is $(W-U^2)$, whose subspace of relations of degree at most two is one-dimensional. Constant linear target changes preserve total degree, and source changes preserve relations. Consequently $\dim_KV_K=1$.

Choose $0\neq R\in V_K$. Its constant term is zero because $H(0)=0$. Over $\Kbar$, its quadratic part has rank one. The rank is also one over $K$, so
\[
 R=\kappa\ell^2+\ell',\qquad \kappa\in K^*,\quad \ell,\ell'\in(K^3)^*.
\]
Here a symmetric rank-one matrix over $K$ can be written as a nonzero scalar times an outer square over $K$ by choosing a nonzero diagonal entry. This does not require a square root. The forms $\ell,\ell'$ are independent. Otherwise $\ell'=c\ell$ and
\[
 \ell(H)\bigl(\kappa\ell(H)+c\bigr)=0.
\]
The polynomial ring is a domain, so $\ell(H)$ is constant. Its value at the origin is zero, contradicting independence of the components. Finally,
\[
 \ell'(H)=-\kappa\ell(H)^2
\]
is the required algebraic dependence over $K$. This completes the proof of Theorem~\ref{thm:dependent-pair}.\qed
